\subsection{Finite Probability, Counting, and Independence}
\begin{definition} \label{def:uniform-probability-countable-space} \hlt{Uniform Probability on a Countable Space}\\
Let $\Omega$ be countable.
\begin{enumerate}[label=(\roman*)]
\setlength{\itemsep}{0pt}
\item A probability measure $\mathbb{P}$ on $(\Omega,\mathcal{P}(\Omega))$ is \hlt{uniform} if $\exists c\in[0,1]$ such that $\mathbb{P}(\{\omega\})=c$ for all $\omega\in\Omega$.
\item The constant $c$ is the \hlt{common atom mass}.
\end{enumerate}
\end{definition}

\begin{proposition} \label{prop:no-uniform-probability-on-countably-infinite-set} \hlt{No Uniform Probability on a Countably Infinite Set}\\
Let $\Omega=\{0,1,2,\ldots\}$. No \hlt{uniform probability measure} exists on $(\Omega,\mathcal{P}(\Omega))$.
\begin{proof}
Assume uniform $\mathbb{P}$ exists, and let $c$ be its common atom mass.\\
Since $\Omega=\bigsqcup\limits_{n=0}^{\infty}\{n\}$, countable additivity gives $1=\mathbb{P}(\Omega)=\sum\limits_{n=0}^{\infty}\mathbb{P}(\{n\})=\sum\limits_{n=0}^{\infty}c$.\\
If $c=0$, the sum is $0$; if $c>0$, the sum diverges. Both contradict $1=\mathbb{P}(\Omega)$.\\
\hlt{Structural Commentary}: Countable additivity forces total mass to equal the sum of the atom masses.\\
\hlt{Key Insight}: Equal mass on countably many atoms is incompatible with total mass $1$.
\end{proof}
\end{proposition}

\begin{method} \label{meth:finite-uniform-probability-by-counting} \hlt{Finite Uniform Probability by Counting}\\
Let $0<|\Omega|<\infty$ and assume the \hlt{uniform model} $\mathbb{P}(\{\omega\})=|\Omega|^{-1}$ for all $\omega\in\Omega$.\\
For $A\subseteq\Omega$, $\mathbb{P}(A)=\sum\limits_{\omega\in A}\mathbb{P}(\{\omega\})=\frac{|A|}{|\Omega|}$, so probability reduces to counting.
\end{method}

\begin{method} \label{meth:finite-uniform-probability-counting-workflow} \hlt{Counting Workflow for Uniform Finite Probability}\\
For a finite uniform sample space $\Omega$ and an event $A\subseteq\Omega$:
\begin{enumerate}[label=(\roman*)]
\setlength{\itemsep}{0pt}
\item choose complete outcomes so that each $\omega\in\Omega$ is equally likely;
\item count $\abs{\Omega}$; and count $\abs{A}$, usually by products, complements, or disjoint cases;
\item conclude $\mathbb{P}(A)=\abs{A}/\abs{\Omega}$.
\end{enumerate}
\hlt{Key Insight}: The hard step is often choosing equally likely outcomes, not applying the ratio.
\end{method}

\begin{example} \label{ex:simpsons-paradox-urn-counting} \hlt{Simpson's Paradox in Two Urn Games}\\
Let $p(U)$ be the red probability in urn $U$, and let $B_i,W_i$ be the black and white urns in game $i$.\\
Within games, $p(B_1)=5/11>3/7=p(W_1)$ and $p(B_2)=6/9>9/14=p(W_2)$.\\
After pooling by color into $B$ and $W$, $p(B)=11/20<12/21=p(W)$, so the aggregate winner is white.\\
\hlt{Key Insight}: A within-group advantage can reverse after aggregation when the group sizes change.
\end{example}

\begin{example} \label{ex:license-plates-and-birthday-complement} \hlt{Complement Counting in Uniform Models}
\begin{enumerate}[label=(\roman*)]
\setlength{\itemsep}{0pt}
\item For plates $\Omega=L^3\times D^3$ with $\abs{L}=26$ and $\abs{D}=10$, $\abs{\Omega}=26^3\cdot10^3$. If $E$ is the event of no repeated letters and no repeated digits, then $\abs{E}=(26)_3(10)_3$, so $\mathbb{P}(E)=\frac{(26)_3(10)_3}{26^3\cdot10^3}$.
\item For $1\leq n\leq365$ iid uniform birthdays, take $\Omega=\{1,\ldots,365\}^n$. If $C$ is the event of a shared birthday, then $C^c=D_n(\{1,\ldots,365\})$, so $\mathbb{P}(C)=1-\frac{(365)_n}{365^n}$; for $n=23$, $\mathbb{P}(C)>1/2$.
\end{enumerate}
\hlt{Key Insight}: ``At least one collision'' is often easier to count by subtracting the no-collision event.
\end{example}

\begin{definition} \label{def:independent-events-and-families} \hlt{Independent Events}
\begin{enumerate}[label=(\roman*)]
\setlength{\itemsep}{0pt}
\item Events $A,B\in\mathcal{F}$ are \hlt{independent}, written $A\perp\!\!\!\perp B$, if $\mathbb{P}(A\cap B)=\mathbb{P}(A)\mathbb{P}(B)$.
\item A family $(A_i)_{i\in I}\subseteq\mathcal{F}$ is \hlt{pairwise independent} if $A_i$ and $A_j$ are independent whenever $i\neq j$.
\item A family $(A_i)_{i\in I}\subseteq\mathcal{F}$ is \hlt{mutually independent} if $\mathbb{P}(\bigcap\limits_{i\in J}A_i)=\prod\limits_{i\in J}\mathbb{P}(A_i)$ for every finite $J\subseteq I$.
\end{enumerate}
\hlt{Key Insight}: Mutual independence means multiplicativity for every finite subfamily, not just pairs.
\end{definition}

\begin{example} \label{ex:pairwise-not-mutual-independence-two-tosses} \hlt{Pairwise Independence Need Not Be Mutual}\\
Let $\Omega=\{HH,HT,TH,TT\}$ with $\mathbb{P}(\{\omega\})=\frac{1}{4}$ $\forall \omega\in\Omega$.\\
Set $A=\{HH,HT\}$, $B=\{HH,TH\}$, $C=\{HT,TH\}$.\\
Then $\mathbb{P}(A)=\mathbb{P}(B)=\mathbb{P}(C)=\frac{1}{2}$, $\mathbb{P}(A\cap B)=\mathbb{P}(A\cap C)=\mathbb{P}(B\cap C)=\frac{1}{4}$, so $A,B,C$ pairwise independent.\\
But $A\cap B\cap C=\emptyset$, so $\mathbb{P}(A\cap B\cap C)=0\neq \frac{1}{8}=\mathbb{P}(A)\mathbb{P}(B)\mathbb{P}(C)$.\\
\hlt{Key Insight}: Pairwise checks do not control higher-order intersections.
\end{example}

\begin{proposition} \label{prop:basic-independence-tests} \hlt{Basic Independence Tests}\\
Let $A,B\in\mathcal{F}$.
\begin{enumerate}[label=(\roman*)]
\setlength{\itemsep}{0pt}
\item If $A\cap B=\emptyset$ and $\mathbb{P}(A)\mathbb{P}(B)>0$, then $A\not\perp\!\!\!\perp B$.
\item If $A\perp\!\!\!\perp B$, then $A^c\perp\!\!\!\perp B$, $A\perp\!\!\!\perp B^c$, and $A^c\perp\!\!\!\perp B^c$.
\item If $\mathbb{P}(A)\in\{0,1\}$, then $A\perp\!\!\!\perp C$ for every $C\in\mathcal{F}$.
\item If $A\perp\!\!\!\perp A$, then $\mathbb{P}(A)\in\{0,1\}$.
\end{enumerate}
\begin{proof}
\begin{enumerate}[label=(\roman*)]
\setlength{\itemsep}{0pt}
\item Since $A\cap B=\emptyset$, $\mathbb{P}(A\cap B)=0\neq\mathbb{P}(A)\mathbb{P}(B)$.
\item If $A\perp\!\!\!\perp B$, then $\mathbb{P}(A^c\cap B)=\mathbb{P}(B)-\mathbb{P}(A\cap B)=(1-\mathbb{P}(A))\mathbb{P}(B)=\mathbb{P}(A^c)\mathbb{P}(B)$\\
and $\mathbb{P}(A\cap B^c)=\mathbb{P}(A)-\mathbb{P}(A\cap B)=\mathbb{P}(A)(1-\mathbb{P}(B))=\mathbb{P}(A)\mathbb{P}(B^c)$.\\
Applying the first identity to $A\perp\!\!\!\perp B^c$ gives $A^c\perp\!\!\!\perp B^c$.
\item Let $C\in\mathcal{F}$. If $\mathbb{P}(A)=0$, then $0\leq\mathbb{P}(A\cap C)\leq\mathbb{P}(A)=0$, so $\mathbb{P}(A\cap C)=0=\mathbb{P}(A)\mathbb{P}(C)$.\\
If $\mathbb{P}(A)=1$, then $\mathbb{P}(A^c)=0$, hence $A^c\perp\!\!\!\perp C$ by the previous case and $A\perp\!\!\!\perp C$ by (ii).
\item Let $p=\mathbb{P}(A)$. From $A\perp\!\!\!\perp A$, $p=\mathbb{P}(A\cap A)=p^2$, hence $p\in\{0,1\}$.
\end{enumerate}
\hlt{Structural Commentary}: Independence is probabilistic, not visual; disjointness usually means strong dependence.\\
\hlt{Key Insight}: Complements preserve independence because $B$ splits into $A\cap B$ and $A^c\cap B$.
\end{proof}
\end{proposition}

\begin{example} \label{ex:independence-complement-first-success-examples} \hlt{Independence, Complements, and First Success}
\begin{enumerate}[label=(\roman*)]
\setlength{\itemsep}{0pt}
\item \hlt{Independence by counting}: for a fair die, $\Omega=\{1,\ldots,6\}$, $A=\{2,4,6\}$, and $B=\{1,2,3,4\}$.\\
Then $\mathbb{P}(A\cap B)=\frac{2}{6}=(\frac{3}{6})(\frac{4}{6})=\mathbb{P}(A)\mathbb{P}(B)$, hence $A\perp\!\!\!\perp B$.
\item \hlt{Complement method}: for independent fair tosses $X_1,\ldots,X_{10}$, let $A=\{\exists j\leq10:X_j=H\}$\\
and $T_j=\{X_j=T\}$. Then $A^c=\bigcap\limits_{j=1}^{10}T_j$ and $\mathbb{P}(A)=1-\mathbb{P}(A^c)=1-\prod\limits_{j=1}^{10}\mathbb{P}(T_j)=1-(\frac{1}{2})^{10}$.
\item \hlt{First-success decomposition}: let $p_1=\frac{1}{3}$, $p_2=\frac{1}{4}$.\\
Let $A_j$ be player $1$'s first success on turn $j$ before player $2$ succeeds.\\
Then $E=\bigsqcup\limits_{j\geq1}A_j$, $\mathbb{P}(A_j)=((1-p_1)(1-p_2))^{j-1}p_1=(\frac{1}{2})^{j-1}/3$, and $\mathbb{P}(E)=\sum\limits_{j\geq1}\mathbb{P}(A_j)=\frac{2}{3}$.
\end{enumerate}
\end{example}

\begin{method} \label{meth:independence-complement-disjoint-decomposition} \hlt{Independence, Complement, and Decomposition Methods}
\begin{enumerate}[label=(\roman*)]
\setlength{\itemsep}{0pt}
\item \hlt{Verify independence}: compute $\mathbb{P}(A)$, $\mathbb{P}(B)$, and $\mathbb{P}(A\cap B)$; conclude $A\perp\!\!\!\perp B$ iff $\mathbb{P}(A\cap B)=\mathbb{P}(A)\mathbb{P}(B)$.
\item \hlt{Use complements}: if $A^c$ is simpler, compute $\mathbb{P}(A)=1-\mathbb{P}(A^c)$.\\
If $A^c=\bigcap\limits_{j=1}^{n}C_j$ with independent $C_j$, then $\mathbb{P}(A)=1-\prod\limits_{j=1}^{n}\mathbb{P}(C_j)$.
\item \hlt{Disjoint decomposition}: write $E=\bigsqcup\limits_jA_j$ with disjoint $A_j$, then $\mathbb{P}(E)=\sum\limits_j\mathbb{P}(A_j)$.
\item \hlt{Geometric summation}: if $\mathbb{P}(A_j)=ar^{j-1}$ for $j\geq1$ and $|r|<1$, then $\mathbb{P}(E)=\sum\limits_{j\geq1}ar^{j-1}=\frac{a}{1-r}$.
\end{enumerate}
\end{method}

\begin{example} \label{ex:reliability-and-match-probabilities} \hlt{Reliability and Match Probabilities}\\
Let $A_i$ be the event that component $i$ works, and assume all relevant events are mutually independent.
\begin{enumerate}[label=(\roman*)]
\setlength{\itemsep}{0pt}
\item Series: if $\mathbb{P}(A_i^c)=p$ for $1\leq i\leq n$ and $W_s=\bigcap\limits_{i=1}^{n}A_i$, then $\mathbb{P}(W_s)=(1-p)^n$ and $\mathbb{P}(W_s^c)=1-(1-p)^n$.
\item Parallel: if $W_p=\bigcup\limits_{i=1}^{n}A_i$, then $\mathbb{P}(W_p^c)=\mathbb{P}(\bigcap\limits_{i=1}^{n}A_i^c)=p^n$.
\item DNA match: let $N_j\in\{A,C,G,T\}$ with $\mathbb{P}(N_j=x)=p_{xj}$ for $j=1,2$.\\
If $N_1\perp\!\!\!\perp N_2$, then $\mathbb{P}(N_1=N_2)=\sum\limits_{x\in\{A,C,G,T\}}p_{x1}p_{x2}$.
\end{enumerate}
\hlt{Key Insight}: Independence multiplies intersections; complements handle failures; matches sum disjoint cases.
\end{example}
